(1) 一个与温度 \(T>0\) 大热库接触的经典二态存储器处于热平衡，状态 \(0,1\) 能量分别为 \(0,\Delta\ge0\)，初始 \(\Delta=0\)。求配分函数 \(Z(\Delta)\)、平衡出错态概率 \(p_1\)、Helmholtz自由能 \(F(\Delta)\)、平均内能 \(U(\Delta)\)。
(2) 外部极缓慢升高 \(\Delta\) 到指定正值，对此可逆等温过程积分实际广义功 \(\delta W=p_1\,d\Delta\)，导出有限终态错误率 \(p_{\rm err}=p_1(\Delta)\) 的最小可逆功。求 \(\Delta\to\infty\) 的Landauer界，注意若不闭合协议，终态有不同能级。
(3) 求初态和有限终态Shannon熵，使用 \(Q=T\Delta S\) 核算环境吸收的热量及 \(\Delta U=W+Q\) 第一律。证明 \(\Delta\to\infty\) 时 \(Q=-k_BT\ln2\)、\(U\to0\)，解释可逆极限与现实有限时间不可逆耗散的区别。
(4) 讨论当允许错误率 \(0<p_{\rm err}<1/2\) 时，为保证错误率至少小于阈值 \(\epsilon\) 应取的最小能级差 \(\Delta_{\min}\)。推导其随 \(\epsilon\to0\) 的对数增长，并区分可逆外源功与最大瞬时储存能量 \(\Delta\)。
